\begin{theorem}[G-20: The exact eventual four-row multiplicity phase]
\label{G-20}
Let $n>m\ge3$, put $a=n-m$, and define
\[
 C_m=\min\left\{\binom{m+3}{3},\,4m^2-2m+2\right\},
 \qquad N_m=\left\lfloor\sqrt{\left\lfloor C_m/2\right\rfloor}\right\rfloor.
\]
For a fixed partition direction $\nu\vdash ns$, $s\ge1$, with
at most four rows, assume the strict conditions
\[
 \nu_4>0,\qquad \nu_1>(n-m)s.
\]
Then exactly one of the following holds:
\begin{enumerate}
\item If $m<n\le N_m$, the padded-permanent multiplicity is
strictly larger than the determinant multiplicity in every
sufficiently large integer stretch $j\nu$.
\item If $n>N_m$, the determinant multiplicity is strictly
larger in every sufficiently large integer stretch $j\nu$.
\end{enumerate}
In both cases the original coordinate degree is $sj$, both
multiplicities are positive for every $j\ge\beta^2$ with
$\beta$ from G-17, and
\begin{equation}\label{G20:exact_ratio}
 \frac{m_{j\nu}(\ell^a\operatorname{per}_m)}
      {m_{j\nu}(\det_n)}
 =\kappa_\nu j^{C_m-2n^2+1}\bigl(1+O(j^{-1})\bigr),
 \qquad \kappa_\nu>0.
\end{equation}
There is no integer equality case in the exponent for the
stated values of $m,n$.

For $m\ge17$, the exact alternatives are
\begin{equation}\label{G20:large_boundary}
 \begin{aligned}
 n^2\le2m^2-m+1&\quad\Longrightarrow\quad
       \text{eventual permanent dominance},\\
 n^2\ge2m^2-m+2&\quad\Longrightarrow\quad
       \text{eventual determinant dominance}.
 \end{aligned}
\end{equation}
The small-size values are
\[
\begin{array}{c|rrrrrrr}
m&3&4&5&6&7&8&9\\ \hline
C_m&20&35&56&84&120&165&220\\
N_m&3&4&5&6&7&9&10
\end{array}
\]
\[
\begin{array}{c|rrrrrrr}
m&10&11&12&13&14&15&16\\ \hline
C_m&286&364&455&560&680&816&969\\
N_m&11&13&15&16&18&20&22.
\end{array}
\]

The corresponding exact criterion for eventual positive
four-row rectangular gaps is
\begin{equation}\label{G20:rectangular}
 m<n\le N_m,\qquad n<4m/3.
\end{equation}
Every admissible rectangular direction then has a positive
gap at every sufficiently large integer stretch. For every
$m\ge14$, condition \eqref{G20:rectangular} simplifies to
$m<n<4m/3$. The threshold $14$ is sharp for this uniform
eventual statement. No assertion is made that an individual
finite-degree gap is impossible in a direction that eventually
has determinant dominance.
\end{theorem}

\begin{proof}
R-13 and R-27 give $c_m=C_m$: sizes through sixteen fill the
four-variable form space, and every size at least seventeen
attains the row--column parameter bound. Substitute this exact
dimension in G-18. Its occurrence theorem gives positivity
of both multiplicities after $\beta^2$, and its ratio formula
becomes \eqref{G20:exact_ratio}.

The exponent is positive precisely when
$2n^2<C_m+1$, equivalently
$n^2\le\lfloor C_m/2\rfloor$. This is precisely $n\le N_m$.
For $m\ge17$, $C_m=4m^2-2m+2$ is even, so the exponent
$C_m-2n^2+1$ is odd and cannot vanish. The two adjacent integer
inequalities are exactly \eqref{G20:large_boundary}.
For $3\le m\le16$, the table is obtained by direct evaluation
of $\binom{m+3}{3}$. The only odd dimensions in that table are
$35,165,455,969$. Their half-successors are respectively
$18,83,228,485$, none of which is a square:
\[
 4^2<18<5^2,\quad9^2<83<10^2,\quad
 15^2<228<16^2,\quad22^2<485<23^2.
\]
Thus the exponent never vanishes in the small-size range
either. Its sign and the positive leading constant in
\eqref{G20:exact_ratio} prove the exact alternatives.

For completeness, the bounds are effective without waiting
for an unspecified asymptotic. With the notation of G-17 set
\[
 u=2n^2+r_\eta-13,\qquad v=C_m+r_\eta-12.
\]
If $v>u$, G-18 gives permanent dominance for every
\[
 j\ge E\bigl(4u(pE)^{u-1}+1\bigr).
\]
If $u>v$, its proof in paragraph~3 applies with the actual
dimension $v$ in place of the upper bound $\bar v$: the
permanent image degree is bounded by $(pE)^{v-1}$, and the
same standard-Veronese and residue-class sandwich gives
determinant dominance for every
\[
 j\ge E\bigl(4v(pE)^{v-1}+1\bigr).
\]
This is the same proved binomial comparison with the exact
dimension substituted, not a new claim about an uncomputed
leading coefficient.

For a rectangle $\nu=(q,q,q,q)$ with $4q=ns$, the first
strict condition is automatic and the second is
$ns/4>(n-m)s$, equivalently $n<4m/3$.
Within that range, the preceding exact sign criterion applies.
At or above its boundary, G-12 gives the opposite weak
multiplicity inequality at every stretch. This proves
\eqref{G20:rectangular}.

R-26 already proves that every $m\ge14$ and every
$m<n<4m/3$ lies in the positive dimension range; its four
small cases and uniform inequality are exact. To prove
sharpness for the eventual statement, take
\[
 m=13,\qquad n=17,\qquad s=4,\qquad\nu=(17,17,17,17).
\]
The direction is strict because $17>4(17-13)=16$.
Here $C_{13}=560$, and the ratio exponent is
\[
 C_{13}-2\cdot17^2+1=560-578+1=-17.
\]
Consequently the ratio is $\kappa_\nu j^{-17}(1+O(j^{-1}))$,
and the determinant multiplicity is larger at every sufficiently
large stretch. Any uniform threshold below fourteen would
include this pair, so it cannot have the asserted eventual
property. This argument does not decide any of its finitely
many earlier stretches.
\end{proof}

\paragraph{Scope.}
The theorem classifies the eventual sign along each fixed
four-row direction. It neither supplies a polynomial-time
positive extraction rule in growing ambient dimension nor
excludes changing-direction or isolated finite-degree gaps.
It does not establish containment of orbit closures, a general
permanent circuit lower bound, or a superlinear-padding GCT
separation. The hypotheses and explicit stretch bounds of
G-17/G-18 remain part of the conclusion.
